% GENERATED FILE. Edit canonical lessons, then run npm run generate:books.
% canonical-source-hash: fnv1a64:7f755454099018ee
% canonical-lessons: FR-C117-bottes, FR-C117-casquette, FR-C117-maillot, FR-C117-taille, FR-C117-portefeuille

\chapter{Boots, Cap, Swimsuit, Size, Wallet}
\label{ch:fr-boots-cap-swimsuit-size-wallet}
\hlchaptermodality{117}

\begin{chapteropening}
\textbf{By the end of this chapter:} \emph{I can name boots, a cap, a swimsuit, size, a wallet.}

Complete the last lesson of chapter 117: I can name boots, a cap, a swimsuit, size, a wallet.
\end{chapteropening}

\section[les bottes]{les bottes — boots}

\label{lesson:FR-C117-bottes}



\begin{quote}
Before the new one: say the French for gloves, then the French for a scarf.
\end{quote}

\subsection*{You'll want to know: les bottes}
\textbf{les bottes} — \textquotedblleft{}boots\textquotedblright{}.

\begin{cousinweb}
\textbf{bottes} is the plural of the word for \textquotedblleft{}boot\textquotedblright{}, whose origin is uncertain. EtymDB 2.1 records Old French \emph{bote}, \textquotedblleft{}a high, thick shoe\textquotedblright{}, and offers a Frankish source that is not secure.
\end{cousinweb}

\begin{grammarlens}[title={What you've built}]
One more word for this chapter.
\end{grammarlens}

\subsection*{Your turn}
\begin{itemize}
\raggedright
  \item \emph{Say it:} \emph{les bottes}
  \item \emph{Say it:} \emph{les bottes}, once more
  \item \emph{Say it:} say \emph{les bottes}
  \item \emph{Recall:} say the French for gloves, then the French for a scarf, then say \emph{les bottes} again
\end{itemize}

\begin{tcolorbox}[breakable,colback=teal!4,colframe=teal!35!black,title={Before you move on}]
What does les bottes mean? (\textquotedblleft{}Boots\textquotedblright{}.) Say it once more.
\end{tcolorbox}
\section[la casquette]{la casquette — a cap}

\label{lesson:FR-C117-casquette}



\begin{quote}
Before the new one: say the French for a scarf, then the French for boots.
\end{quote}

\subsection*{You'll want to know: la casquette}
\textbf{la casquette} — \textquotedblleft{}a cap\textquotedblright{}.

\begin{cousinweb}
\textbf{casquette} — a little form of the French word for \textquotedblleft{}helmet\textquotedblright{}, with \emph{-ette}. That word came from Spanish \emph{casco}, \textquotedblleft{}skull, helmet\textquotedblright{} (EtymDB 2.1).
\end{cousinweb}

\begin{grammarlens}[title={What you've built}]
One more word for this chapter.
\end{grammarlens}

\subsection*{Your turn}
\begin{itemize}
\raggedright
  \item \emph{Say it:} \emph{la casquette}
  \item \emph{Say it:} \emph{la casquette}, once more
  \item \emph{Say it:} say \emph{la casquette}
  \item \emph{Recall:} say the French for a scarf, then the French for boots, then say \emph{la casquette} again
\end{itemize}

\begin{tcolorbox}[breakable,colback=teal!4,colframe=teal!35!black,title={Before you move on}]
What does la casquette mean? (\textquotedblleft{}A cap\textquotedblright{}.) Say it once more.
\end{tcolorbox}
\section[le maillot]{le maillot — a swimsuit; a sports shirt}

\label{lesson:FR-C117-maillot}



\begin{quote}
Before the new one: say the French for boots, then the French for a cap.
\end{quote}

\subsection*{You'll want to know: le maillot}
\textbf{le maillot} — \textquotedblleft{}a swimsuit; a sports shirt\textquotedblright{}.

\begin{cousinweb}
\textbf{maillot} $\leftarrow$ Old French \emph{mailloel}, a swaddling band, from \emph{maille}, \textquotedblleft{}a stitch, a mesh\textquotedblright{}, $\leftarrow$ Latin \emph{macula}, \textquotedblleft{}a spot, a mesh\textquotedblright{}. (EtymDB 2.1)
\end{cousinweb}

\begin{grammarlens}[title={What you've built}]
One more word for this chapter.
\end{grammarlens}

\subsection*{Your turn}
\begin{itemize}
\raggedright
  \item \emph{Say it:} \emph{le maillot}
  \item \emph{Say it:} \emph{le maillot}, once more
  \item \emph{Say it:} say \emph{le maillot}
  \item \emph{Recall:} say the French for boots, then the French for a cap, then say \emph{le maillot} again
\end{itemize}

\begin{tcolorbox}[breakable,colback=teal!4,colframe=teal!35!black,title={Before you move on}]
What does le maillot mean? (\textquotedblleft{}A swimsuit; a sports shirt\textquotedblright{}.) Say it once more.
\end{tcolorbox}
\section[la taille]{la taille — size; the waist}

\label{lesson:FR-C117-taille}



\begin{quote}
Before the new one: say the French for a cap, then the French for a swimsuit; a sports shirt.
\end{quote}

\subsection*{You'll want to know: la taille}
\textbf{la taille} — \textquotedblleft{}size; the waist\textquotedblright{}.

\begin{cousinweb}
\textbf{taille} — formed from the verb for \textquotedblleft{}to cut\textquotedblright{}, $\leftarrow$ Late Latin \emph{taliare}, from \emph{talea}, \textquotedblleft{}a cutting, a twig\textquotedblright{} (EtymDB 2.1): the cut of a garment, then its size.
\end{cousinweb}

\begin{grammarlens}[title={What you've built}]
One more word for this chapter.
\end{grammarlens}

\subsection*{Your turn}
\begin{itemize}
\raggedright
  \item \emph{Say it:} \emph{la taille}
  \item \emph{Say it:} \emph{la taille}, once more
  \item \emph{Say it:} say \emph{la taille}
  \item \emph{Recall:} say the French for a cap, then the French for a swimsuit; a sports shirt, then say \emph{la taille} again
\end{itemize}

\begin{tcolorbox}[breakable,colback=teal!4,colframe=teal!35!black,title={Before you move on}]
What does la taille mean? (\textquotedblleft{}Size; the waist\textquotedblright{}.) Say it once more.
\end{tcolorbox}
\section[le portefeuille]{le portefeuille — a wallet}

\label{lesson:FR-C117-portefeuille}



\begin{quote}
Before the new one: say the French for a swimsuit; a sports shirt, then the French for size; the waist.
\end{quote}

\subsection*{You'll want to know: le portefeuille}
\textbf{le portefeuille} — \textquotedblleft{}a wallet\textquotedblright{}.

\begin{cousinweb}
\textbf{portefeuille} — \textquotedblleft{}carry-leaves\textquotedblright{}: \emph{porte}, from \emph{porter}, \textquotedblleft{}to carry\textquotedblright{}, and \emph{feuille}, \textquotedblleft{}a leaf, a sheet\textquotedblright{}: a case for carrying papers.
\end{cousinweb}

\begin{grammarlens}[title={What you've built}]
One more word for this chapter.
\end{grammarlens}

\subsection*{Your turn}
\begin{itemize}
\raggedright
  \item \emph{Say it:} \emph{le portefeuille}
  \item \emph{Say it:} \emph{le portefeuille}, once more
  \item \emph{Say it:} say \emph{le portefeuille}
  \item \emph{Recall:} say the French for a swimsuit; a sports shirt, then the French for size; the waist, then say \emph{le portefeuille} again
\end{itemize}

\begin{tcolorbox}[breakable,colback=teal!4,colframe=teal!35!black,title={Before you move on}]
What does le portefeuille mean? (\textquotedblleft{}A wallet\textquotedblright{}.) Say it once more.
\end{tcolorbox}
